% !TeX root = ../main.tex
\section{Design decisions}
\label{sec:design}

Building a DSL is a sequence of decisions about syntax, grammar, types,
semantics, execution and errors, and the decisions interact. A more
expressive language is harder to implement, analyze and maintain; a more
restricted one leaves users without a way to say some things. This section
goes beyond the slides, and uses the training DSL of
Section~\ref{sec:dsl} as its example.

\subsection{Syntax}

\paragraph{Who writes the programs.} Programmers accept braces, semicolons and
type annotations in exchange for a familiar structure. Domain experts need a
notation that reads like their domain. The users of the training DSL are
coaches and athletes, so its syntax follows the way they describe a session.

\paragraph{Declarative and imperative.} A declarative program says what
exists; an imperative program says what to do. The training plan
\texttt{Monday: run: 10km @threshold} describes a session. The internal Python
version of Section~\ref{sec:dsl} gives the same information as a sequence of
calls, \texttt{.on('Monday').run(10, intensity='threshold')}, that build the
plan step by step. The declarative version talks about the session, not about
the operations that record it.

\paragraph{Verbosity.} A verbose syntax is explicit and easy to read the first
time:
\begin{lstlisting}
exercise type=run distance=10 unit=km intensity=threshold
\end{lstlisting}
A concise syntax has to be learned but is faster to write:
\begin{lstlisting}
run: 10km @threshold
\end{lstlisting}
If programs are written once and read many times, readability matters more;
if the same users write programs every day, conciseness does. The training
DSL is concise.

\paragraph{Layout.} When indentation has meaning, as in Python, the lexer
has to produce tokens for the start and end of each indented block. When
blocks are delimited explicitly, as in C, the lexer can discard all
whitespace. The Rascal grammar of the training DSL takes the second option:
its layout discards spaces and newlines, and day names start new blocks, so
the indentation of the example is only a convention for the reader.

\subsection{Grammar}

\paragraph{Ambiguity.} A grammar is ambiguous when a word has two parse trees
(formal languages notes). In a DSL an ambiguity means that a program can mean
two things. If the grammar allowed a second qualifier after the intensity, in
\begin{lstlisting}
run 10 km @threshold easy
\end{lstlisting}
the word \texttt{easy} could modify the intensity or the whole activity. The
grammar has to decide, with separators, with different tokens for different
qualifiers, or with rules that allow only one reading.

\paragraph{Context-free is enough.} The syntax of most DSLs is context-free,
and with an LL(1) or LR(1) grammar it is parsed in linear time.
Context-sensitive conditions, such as the unit that an exercise accepts, do
not belong in the grammar but in semantic analysis, after the tree is built.

\paragraph{LL or LR.} An LL(1) grammar can be parsed by recursive descent, one
function per nonterminal, and the error messages can be precise; it has to be
left-factored and without left recursion. LR grammars cover more languages,
but their tables are built by generators such as Yacc, Bison or the parser
generator of Rascal. For a small DSL written by hand, LL(1) is the usual
choice (parsing notes). The grammar of the training DSL is LL(1).

\paragraph{Error recovery.} A parser that stops at the first error makes the
user fix one error at a time. Panic mode, from the parsing lecture, skips
tokens until a synchronization token (a new day name, in the training DSL)
and continues from there. The message is part of the design:
\begin{lstlisting}
line 2: invalid unit 'min' for exercise 'run'
\end{lstlisting}
says more than \texttt{parse error at token 'min'}.

\subsection{Architecture of the interpreter}

The decisions meet in the implementation (Figure~\ref{fig:architecture}). Each
component has one responsibility: the lexer does not know the grammar, the
parser does not know the rules of the domain, and the checker does not know
how the plan is executed. Each component can be tested on its own, and every
error is reported by the component that found it.

\begin{figure}[h]
\centering
\begin{tikzpicture}[st/.style={draw, minimum width=3.2cm, minimum height=7mm, font=\small}, er/.style={draw, dashed, font=\footnotesize, align=center}, node distance=5mm]
  \node[st] (file) {\texttt{plan.txt}};
  \node[st, below=of file] (lex) {lexer};
  \node[st, below=of lex] (par) {parser (LL(1))};
  \node[st, below=of par] (check) {semantic checker};
  \node[st, below=of check] (ev) {evaluator};
  \node[st, below=of ev] (out) {training schedule};
  \draw[flecha] (file) -- (lex);
  \draw[flecha] (lex) -- node[right, font=\scriptsize] {tokens} (par);
  \draw[flecha] (par) -- node[right, font=\scriptsize] {AST} (check);
  \draw[flecha] (check) -- node[right, font=\scriptsize] {checked AST} (ev);
  \draw[flecha] (ev) -- (out);
  \node[er, right=2.2cm of lex] (e1) {lexical error\\ (bad character)};
  \node[er, right=2.2cm of par] (e2) {syntax error\\ (bad structure)};
  \node[er, right=2.2cm of check] (e3) {semantic error\\ (bad meaning)};
  \draw[flecha, dashed] (lex) -- (e1);
  \draw[flecha, dashed] (par) -- (e2);
  \draw[flecha, dashed] (check) -- (e3);
\end{tikzpicture}
\caption{Architecture of an interpreter for the training DSL.}
\label{fig:architecture}
\end{figure}

\subsection{Abstraction level}

A DSL should use the abstractions of its domain, no higher and no lower
(Figure~\ref{fig:abstraction}). A language at the level of a GPL makes the
user think as a programmer. Natural language loses precision and needs a
translation back into a structured form. When a DSL makes its user think about
the level below, the abstraction leaks: an error message such as
\texttt{NullPointerException in ActivityBuilder.java:43} speaks about the
implementation, which a coach does not know.

\begin{figure}[h]
\centering
\begin{tikzpicture}[lv/.style={draw, minimum width=8cm, minimum height=7mm, font=\small}, node distance=2mm]
  \node[lv] (nat) {natural language (``run 10k easy tomorrow'')};
  \node[lv, below=of nat, very thick] (dsl) {domain abstractions (plan, session, intensity)};
  \node[lv, below=of dsl] (gpl) {GPL abstractions (classes, modules, types)};
  \node[lv, below=of gpl] (imp) {imperative code (loops, data structures, calls)};
  \node[lv, below=of imp] (mach) {machine code or bytecode};
  \node[font=\small, right=3mm of dsl] {the DSL};
  \draw[flecha] ([xshift=-4.6cm]mach.south) -- node[left, font=\scriptsize, align=right] {more\\ abstract} ([xshift=-4.6cm]nat.north);
\end{tikzpicture}
\caption{Levels of abstraction. The DSL is at the level of the domain.}
\label{fig:abstraction}
\end{figure}

\subsection{Types and domain constraints}

Programs that respect the grammar can still be meaningless:
\texttt{run: 10min} gives a duration to an exercise measured in distance, and
\texttt{rest: 5km} the opposite. A type system for the DSL rejects them before
the program runs. The type of each exercise is known from its name:
\texttt{run} has type $\mathit{Distance} \to \mathit{Activity}$, and
\texttt{rest} has type $\mathit{Duration} \to \mathit{Activity}$, so the user
never writes a type annotation. The checker of the training DSL implements
these rules, and also requires positive amounts:
\coderange{Rascal}{rascal/src/Training.rsc}{check}
The tests check the messages for \texttt{run: 10min}, \texttt{rest: 5km} and
\texttt{run: 0km}. A negative amount, \texttt{run: -5km}, does not even reach
the checker: the grammar has no minus sign, so it is a syntax error.

Constraints of this kind can be written as types with a predicate, called
refinement types:
\[
  \mathit{Distance} = \{ d \in \mathbb{N} \mid d > 0 \}.
\]
The constraints are checked statically, when the program is analyzed. A plan
is stored and used weeks after it is written, and a dynamic check would only
find the error when the plan runs. The rule for the training DSL is to reject
a wrong plan before it is stored.

\subsection{Semantics and execution}

Section~\ref{sec:dsl} gave an operational and a denotational rule for one
activity. For a whole plan, the denotation of a plan is the set of
denotations of its days,
\[
  \llbracket \mathit{plan} \rrbracket = \{ \llbracket d \rrbracket \mid d \in \mathit{days}(\mathit{plan}) \},
\]
and two programs with the same denotation mean the same, even if they are
written differently. A simple evaluator walks the checked tree and computes,
for instance, the total distance of the plan:
\coderange{Rascal}{rascal/src/Training.rsc}{eval}
For the example plan the result is 12000 meters, 10 km of running and 2000 m
of swimming.

The execution model can be an interpreter, a compiler or a translator to
another high-level language (Figure~\ref{fig:execution}). For a DSL an
interpreter is the usual first choice: there is no build step, and errors
point to the source. Compiling pays off when the domain involves heavy
computation, and translating to another language when the programs must run
in an existing ecosystem, as TypeScript translated to JavaScript. The training
DSL produces schedules and totals, and an interpreter is enough.

\begin{figure}[h]
\centering
\begin{tikzpicture}[st/.style={draw, minimum width=2.6cm, minimum height=7mm, font=\small}]
  \node[st] (ast) at (0,0) {checked AST};
  \node[st] (i) at (4,1.1) {interpreter};
  \node[st] (c) at (4,0) {compiler};
  \node[st] (t) at (4,-1.1) {translator};
  \node[font=\small, anchor=west] (ri) at (5.6,1.1) {result};
  \node[font=\small, anchor=west] (rc) at (5.6,0) {bytecode or machine code};
  \node[font=\small, anchor=west] (rt) at (5.6,-1.1) {source of another language};
  \draw[flecha] (ast.east) -- (i.west);
  \draw[flecha] (ast.east) -- (c.west);
  \draw[flecha] (ast.east) -- (t.west);
  \draw[flecha] (i) -- (ri); \draw[flecha] (c) -- (rc); \draw[flecha] (t) -- (rt);
\end{tikzpicture}
\caption{Three execution models for a DSL.}
\label{fig:execution}
\end{figure}

\subsection{Errors}

Errors in a DSL come in four levels, each found by a different check:
\begin{enumerate}
  \item lexical: \texttt{r\%n: 10km} has a character outside the alphabet;
  \item syntactic: \texttt{run 10km:} has valid tokens in the wrong order;
  \item semantic: \texttt{run: 10min} has a valid structure and the wrong unit;
  \item domain: \texttt{run: 500km @sprint} is valid and well typed, but no
    session can have it.
\end{enumerate}
A message should say what is wrong, where, and what was expected, in the terms
of the domain.

\subsection{A natural language layer}

The training DSL was designed to accept natural language input, such as
``add a threshold run on Monday, about 10k'', translated into a DSL program by
a language model before the lexer (Figure~\ref{fig:nl-layer}). The layer is
not part of the DSL. The DSL acts as a structured intermediate representation:
the model does not produce schedules but DSL programs, which go through the
whole pipeline. If the program has a wrong unit, an impossible value or a
broken structure, the checker rejects it, and its message can be sent back to
the model to ask for a correction. The model deals with the ambiguity of
natural language, and the DSL with structure and correctness.

\begin{figure}[h]
\centering
\begin{tikzpicture}[st/.style={draw, minimum width=3.6cm, minimum height=7mm, font=\small}, node distance=6mm]
  \node[st] (nl) {natural language input};
  \node[st, below=of nl] (ai) {translation layer};
  \node[st, below=of ai] (dsl) {DSL program};
  \node[st, below=of dsl] (pipe) {DSL pipeline};
  \node[st, below=of pipe] (out) {training schedule};
  \draw[flecha] (nl) -- node[right, font=\scriptsize] {free text} (ai);
  \draw[flecha] (ai) -- (dsl);
  \draw[flecha] (dsl) -- (pipe);
  \draw[flecha] (pipe) -- node[right, font=\scriptsize] {checked} (out);
  \node[draw, dashed, font=\footnotesize, align=center, right=2cm of dsl] (fb) {error messages};
  \draw[flecha, dashed] (pipe.east) -| (fb.south);
  \draw[flecha, dashed] (fb.north) |- (ai.east);
\end{tikzpicture}
\caption{A natural language layer in front of the DSL.}
\label{fig:nl-layer}
\end{figure}

\subsection{Evolution}

A DSL whose changes break the programs already written is a liability, and
stored plans are exactly such programs. New terminals, such as a new exercise,
and new optional constructs do not affect programs that do not use them. A
new required construct breaks old programs, and needs a version: a header such
as \texttt{version: 2} tells the parser which grammar to use.

Table~\ref{tab:decisions} collects the decisions taken for the training DSL.

\begin{table}[h]
\centering
\begin{tabular}{@{}p{3.4cm}p{3.4cm}p{6cm}@{}}
\toprule
Decision & Choice & Reason \\
\midrule
Style & declarative & coaches describe sessions \\
Syntax & concise & the same users write plans every day \\
Layout & ignored by the grammar & simpler lexer; day names start blocks \\
Parser & LL(1) & small grammar, precise errors \\
Types & inferred, static & no annotations, errors before storing \\
Execution & interpreted & no performance requirement \\
Errors & in terms of the domain & users are not programmers \\
Natural language & separate layer & the semantics of the DSL stays simple \\
Evolution & additive changes & stored plans keep working \\
\bottomrule
\end{tabular}
\caption{Design decisions for the training DSL.}
\label{tab:decisions}
\end{table}

\subsection{Exercises}

\begin{exercise}
Add \texttt{cycle: 2h} to a plan. Does the checker accept it? Which change to
the checker would accept both distances and durations for \texttt{cycle}?
\end{exercise}

\paragraph{Solution.} The checker rejects it: \texttt{cycle} is in
\texttt{distanceExercises}, so it expects \texttt{km} or \texttt{m}, and the
message is \texttt{invalid unit 'h' for exercise 'cycle'}. Accepting both
needs a third set, for exercises that accept any unit, and for those the
expected units are the union of the two sets. The total distance would then
have to skip the activities of \texttt{cycle} measured in hours.

\begin{exercise}
Classify each error: \texttt{Monday: run: 10km @Threshold},
\texttt{Monday: run: 10min}, \texttt{Monday run: 10km}, \texttt{Monday: run: 90km}.
\end{exercise}

\paragraph{Solution.} \texttt{@Threshold} is a syntax error: \texttt{Label} only
has lowercase letters, so \texttt{T} cannot start the label, and Rascal reports
the position. \texttt{run: 10min} is a semantic error (wrong unit). The missing
colon after \texttt{Monday} is a syntax error. \texttt{run: 90km} passes every
check of the implementation; it is a domain error that a rule on the maximum
distance of a session would catch.

\begin{exercise}
The example of the Objectilang statement declares \texttt{List[+A]} and
assigns a value of type \texttt{List[Dog]} to a value declared with type
\texttt{List[Animal]}. Which phase checks the assignment, and why is it valid?
\end{exercise}

\paragraph{Solution.} The type checker, in the semantic analysis. With
\texttt{Dog} $<:$ \texttt{Animal} and \texttt{A} covariant,
\texttt{List[Dog]} $<:$ \texttt{List[Animal]}, and by subsumption a value of the
subtype can be used where the supertype is expected (subtyping section of the
types notes).
